%%%%%%%%%%%%%%%%%%%%%%%%%%%%%%%%%%%%%%%%%%%%%%%%%%%%%%%%%%%%%%%%%%%%%%%%%%%%%%%%

\part{Quantum walk}
\begin{frame}
  \partpage
\label{part:walk}
\end{frame}

%%%%%%%%%%%%%%%%%%%%%%%%%%%%%%%%%%%%%%%%%%%%%%%%%%%%%%%%%%%%%%%%%%%%%%%%%%%%%%%%

\begin{frame}{Randomized algorithms}
  
Randomness is an important tool in computer science

\pause
\bigskip

Black-box problems \pause
\begin{itemize}[<+->]
  \item Huge speedups are possible (Deutsch-Jozsa: $2^{\Omega(n)}$ vs.\ $O(1)$)
  \item Polynomial speedup for some total functions (game trees: $\Omega(n)$ vs.\ $O(n^{0.754})$)
\end{itemize}

\pause
\bigskip

Natural problems \pause
\begin{itemize}[<+->]
  \item Majority view is that derandomization should be possible (P=BPP)
  \item Randomness may give polynomial speedups (Sch\"oning algorithm for $k$-SAT)
  \item Can be useful for algorithm design
\end{itemize}

\end{frame}

%%%%%%%%%%%%%%%%%%%%%%%%%%%%%%%%%%%%%%%%%%%%%%%%%%%%%%%%%%%%%%%%%%%%%%%%%%%%%%%%

\begin{frame}{Random walk}

Graph $G=(V,E)$ \hspace{3ex}
\raisebox{-1cm}{
\setlength{\unitlength}{.2cm}
\begin{picture}(18.5,11)(-.5,-.5)
%  \put(-.5,-.5){\framebox(18.5,11){}}
  \put(0,0){\circle*{1}}
  \put(0,10){\circle*{1}}
  \put(10,0){\circle*{1}}
  \put(10,10){\circle*{1}}
  \put(17.5,5){\circle*{1}}
  \thicklines
  \put(0,0){\line(1,0){10}}
  \put(0,0){\line(0,1){10}}
  \put(0,10){\line(1,0){10}}
  \put(10,0){\line(0,1){10}}
  \put(10,0){\line(3,2){7.5}}
  \put(10,10){\line(3,-2){7.5}}
\end{picture}
}

\pause
\bigskip
\bigskip

Two kinds of walks:
\begin{itemize}
  \item Discrete time
  \item Continuous time
\end{itemize}

\end{frame}

%%%%%%%%%%%%%%%%%%%%%%%%%%%%%%%%%%%%%%%%%%%%%%%%%%%%%%%%%%%%%%%%%%%%%%%%%%%%%%%%

\begin{frame}{Random walk algorithms}

Undirected $s$--$t$ connectivity in log space
\begin{itemize}
  \item Problem: given an undirected graph $G=(V,E)$ and $s,t \in V$, is there a path from $s$ to $t$?
  \item A random walk from $s$ eventually reaches $t$ iff there is a path
  \item Taking a random walk only requires log space
  \item Can be derandomized (Reingold 2004), but this is nontrivial
\end{itemize}

\pause
\bigskip

Markov chain Monte Carlo
\begin{itemize}
  \item Problem: sample from some probability distribution (uniform distribution over some set of combinatorial objects, thermal equilibrium state of a physical system, etc.)
  \item Create a Markov chain whose stationary distribution is the desired one
  \item Run the chain until it converges
\end{itemize}

\end{frame}

%%%%%%%%%%%%%%%%%%%%%%%%%%%%%%%%%%%%%%%%%%%%%%%%%%%%%%%%%%%%%%%%%%%%%%%%%%%%%%%%

\begin{frame}{Continuous-time quantum walk}

\begin{columns}
\begin{column}{2cm}
\noindent Graph $G$ \medskip
\setlength{\unitlength}{.12cm}
\begin{picture}(19.5,13)(-.5,-.5)
%  \put(-.5,-.5){\framebox(19.5,13){}}
  \put(0,0){\circle*{1}}
  \put(0,10){\circle*{1}}
  \put(10,0){\circle*{1}}
  \put(10,10){\circle*{1}}
  \put(17.5,5){\circle*{1}}
  \thicklines
  \put(0,0){\line(1,0){10}}
  \put(0,0){\line(0,1){10}}
  \put(0,10){\line(1,0){10}}
  \put(10,0){\line(0,1){10}}
  \put(10,0){\line(3,2){7.5}}
  \put(10,10){\line(3,-2){7.5}}
  \footnotesize
  \put(0,10){\makebox(2,3){1}}
  \put(8,10){\makebox(2,3){2}}
  \put(0,0){\makebox(2,3){3}}
  \put(8,0){\makebox(2,3){4}}
  \put(17.5,5){\makebox(2,3){5}}
\end{picture}
\end{column}
\begin{column}{3cm}
\[
\setlength{\arraycolsep}{.75ex}
A = \footnotesize\begin{pmatrix}
0 & 1 & 1 & 0 & 0 \\
1 & 0 & 0 & 1 & 1 \\
1 & 0 & 0 & 1 & 0 \\
0 & 1 & 1 & 0 & 1 \\
0 & 1 & 0 & 1 & 0
\end{pmatrix}
\]
\begin{center}
{\color{teal} adjacency matrix}
\end{center}
\end{column}
\begin{column}{4cm}
\onslide<2->
\[
\setlength{\arraycolsep}{.5ex}
L = \footnotesize\begin{pmatrix}
-2 & 1 & 1 & 0 & 0 \\
1 & -3 & 0 & 1 & 1 \\
1 & 0 & -2 & 1 & 0 \\
0 & 1 & 1 & -3 & 1 \\
0 & 1 & 0 & 1 & -2
\end{pmatrix}
\]
\begin{center}
{\color{teal} Laplacian}
\end{center}
\end{column}
\end{columns}

\pause \pause
\bigskip

Random walk on $G$
\begin{itemize}
  \item State: probability $p_v(t)$ of being at vertex $v$ at time $t$ \pause
  \item Dynamics: $\frac{\d}{\d{t}}\vec p(t) = -L\vec p(t)$
\end{itemize}

\pause
\bigskip

Quantum walk on $G$
\begin{itemize}
  \item State: amplitude $q_v(t)$ to be at vertex $v$ at time $t$ \\
  (i.e., $|\psi(t)\> = \sum_{v \in V} q_v(t) |v\>$) \pause
  \item Dynamics: $\ii\frac{\d}{\d{t}}\vec q(t) = -L\vec q(t)$
\end{itemize}

\end{frame}

%%%%%%%%%%%%%%%%%%%%%%%%%%%%%%%%%%%%%%%%%%%%%%%%%%%%%%%%%%%%%%%%%%%%%%%%%%%%%%%%

\begin{frame}{Random vs.\ quantum walk on the line}

\begin{center}
\setlength{\unitlength}{.1cm}
\begin{picture}(100,5)(-10,-4)
%  \put(-10,-4){\framebox(100,5){}}
  \multiput(0,0)(10,0){9}{\circle*{1}}
  \thicklines
  \put(0,0){\line(1,0){80}}
  \put(0,0){\vector(-1,0){10}}
  \put(80,0){\vector(1,0){10}}
  \put(-1,-4){\makebox(2,3){-4}}
  \put(9,-4){\makebox(2,3){-3}}
  \put(19,-4){\makebox(2,3){-2}}
  \put(29,-4){\makebox(2,3){-1}}
  \put(39,-4){\makebox(2,3){0}}
  \put(49,-4){\makebox(2,3){1}}
  \put(59,-4){\makebox(2,3){2}}
  \put(69,-4){\makebox(2,3){3}}
  \put(79,-4){\makebox(2,3){4}}
\end{picture}
\end{center}

\pause
\bigskip

\begin{tabular}{cc}
\raisebox{1.5cm}{Classical:} &
  \includegraphics[width=7cm]{linec}
\\[.5cm] \pause
\raisebox{1.5cm}{Quantum:} &
  \includegraphics[width=7cm]{lineq}
\end{tabular}

\end{frame}

%%%%%%%%%%%%%%%%%%%%%%%%%%%%%%%%%%%%%%%%%%%%%%%%%%%%%%%%%%%%%%%%%%%%%%%%%%%%%%%%

\begin{frame}{Random vs.\ quantum walk on the hypercube}

\begin{columns}
\begin{column}{5cm}
\begin{align*}
V = \{&0,1\}^n \\[2pt]
E = \{&(x,y) \in V \times V: \\[-2pt]
&\text{$x$ and $y$ differ in} \\[-2pt]
&\text{exactly one bit}\}
\end{align*}
\end{column}

\begin{column}{5cm}
\raisebox{2.75cm}{$n=3$:}
\setlength{\unitlength}{.15cm}
\begin{picture}(21,21)(-3,-3)
%  \put(-3,-3){\framebox(21,21){}}
  \put(0,0){\circle*{1}}
  \put(0,10){\circle*{1}}
  \put(10,0){\circle*{1}}
  \put(10,10){\circle*{1}}
  \put(5,5){\circle*{1}}
  \put(5,15){\circle*{1}}
  \put(15,5){\circle*{1}}
  \put(15,15){\circle*{1}}
  \thicklines
  \put(0,0){\line(1,0){10}}
  \put(0,0){\line(0,1){10}}
  \put(0,0){\line(1,1){5}}
  \put(10,10){\line(-1,0){10}}
  \put(10,10){\line(0,-1){10}}
  \put(10,10){\line(1,1){5}}
  \put(5,15){\line(1,0){10}}
  \put(5,15){\line(0,-1){10}}
  \put(5,15){\line(-1,-1){5}}
  \put(15,5){\line(-1,0){10}}
  \put(15,5){\line(0,1){10}}
  \put(15,5){\line(-1,-1){5}}
  \footnotesize
  \put(-2,-3){\makebox(4,2){000}}
  \put(8,-3){\makebox(4,2){100}}
  \put(4,2){\makebox(4,2){010}}
  \put(-3,11){\makebox(4,2){001}}
  \put(14,2){\makebox(4,2){110}}
  \put(3,16){\makebox(4,2){011}}
  \put(7,11){\makebox(4,2){101}}
  \put(13,16){\makebox(4,2){111}}
\end{picture}
\end{column}
\end{columns}

\pause
\bigskip

Classical random walk: reaching $11\ldots1$ from $00\ldots0$ is exponentially unlikely

\bigskip

\pause
Quantum walk: with $A = \sum_{j=1}^n X_j$, \pause
\begin{align*}
  e^{-\ii A t}
  = \prod_{j=1}^n e^{-\ii X_j t}
  \onslide<5->{
  = \bigotimes_{j=1}^n \begin{pmatrix} \cos t & -\ii\sin t \\ -\ii\sin t & \cos t \end{pmatrix}}
\end{align*}

\end{frame}

%%%%%%%%%%%%%%%%%%%%%%%%%%%%%%%%%%%%%%%%%%%%%%%%%%%%%%%%%%%%%%%%%%%%%%%%%%%%%%%%

\begin{frame}{Glued trees problem}

\begin{center}
\includegraphics[width=6cm]{gluedtrees}
\end{center}

Black-box description of a graph
\begin{itemize}
  \item Vertices have arbitrary labels
  \item Label of `\textrm{in}' vertex is known
  \item Given a vertex label, black box returns labels of its neighbors
  \item Restricts algorithms to explore the graph locally
\end{itemize}

\end{frame}

%%%%%%%%%%%%%%%%%%%%%%%%%%%%%%%%%%%%%%%%%%%%%%%%%%%%%%%%%%%%%%%%%%%%%%%%%%%%%%%%

\begin{frame}{Glued trees problem: Classical query complexity}

\begin{center}
\includegraphics[width=6cm]{gluedtrees}
\end{center}

Let $n$ denote the height of one of the binary trees

\pause
\bigskip

Classical random walk from `\textrm{in}': probability of reaching `\textrm{out}' is $2^{-\Omega(n)}$ at all times

\pause
\bigskip

In fact, the classical query complexity is $2^{\Omega(n)}$


\end{frame}

%%%%%%%%%%%%%%%%%%%%%%%%%%%%%%%%%%%%%%%%%%%%%%%%%%%%%%%%%%%%%%%%%%%%%%%%%%%%%%%%

\begin{frame}{Glued trees problem: Exponential speedup}

\begin{columns}
\begin{column}{6cm}
\begin{center}
\includegraphics[width=5cm]{gluedtrees}

\bigskip

\onslide<3>{
$\downarrow$

\bigskip

\includegraphics[width=5cm]{gluedtrees-reduced}
}
\end{center}
\end{column}

\begin{column}{6cm}
\onslide<2->{
  Column subspace
  \begin{align*}
  &|\text{col $j$}\> := \frac{1}{\sqrt{N_j}} \sum_{v \in \text{column $j$}}|v\>
  \\
  &N_j := \begin{cases}
  2^j & \text{if $j \in [0,n]$} \\
  2^{2n+1-j} & \text{if $j \in [n+1,2n+1]$}
  \end{cases}
  \end{align*}
}
\onslide<3>{
  \bigskip
  
  Reduced adjacency matrix
  \begin{align*}
  &\<\text{col $j$}|A|\text{col $j+1$}\> \\
  &\quad= \small\begin{cases}
  \sqrt{2} & \text{if $j \in [0,n-1]$} \\
  \sqrt{2} & \text{if $j \in [n+1,2n]$} \\
  2 & \text{if $j=n$}
  \end{cases}
  \end{align*}
}
\end{column}
\end{columns}

\end{frame}

%%%%%%%%%%%%%%%%%%%%%%%%%%%%%%%%%%%%%%%%%%%%%%%%%%%%%%%%%%%%%%%%%%%%%%%%%%%%%%%%

\begin{frame}{Discrete-time quantum walk: Need for a coin}

Quantum analog of discrete-time random walk?

\pause
\bigskip

Unitary matrix $U \in \C^{|V| \times |V|}$ with $U_{vw} \ne 0$ iff $(v,w) \in E$

\pause
\bigskip

Consider the line:

\begin{center}
\setlength{\unitlength}{.1cm}
\begin{picture}(100,5)(-10,-4)
%  \put(-10,-4){\framebox(100,5){}}
  \multiput(0,0)(10,0){9}{\circle*{1}}
  \thicklines
  \put(0,0){\line(1,0){80}}
  \put(0,0){\vector(-1,0){10}}
  \put(80,0){\vector(1,0){10}}
  \put(-1,-4){\makebox(2,3){-4}}
  \put(9,-4){\makebox(2,3){-3}}
  \put(19,-4){\makebox(2,3){-2}}
  \put(29,-4){\makebox(2,3){-1}}
  \put(39,-4){\makebox(2,3){0}}
  \put(49,-4){\makebox(2,3){1}}
  \put(59,-4){\makebox(2,3){2}}
  \put(69,-4){\makebox(2,3){3}}
  \put(79,-4){\makebox(2,3){4}}
\end{picture}
\end{center}

\pause

Define walk by $|x\> \mapsto \frac{1}{\sqrt2}(|x-1\> + |x+1\>)$?

\pause
\bigskip

But then $|x+2\> \mapsto \frac{1}{\sqrt2}(|x+1\>+|x+3\>)$,
so this is not unitary!

\pause
\bigskip

In general, we must enlarge the state space.

\end{frame}

%%%%%%%%%%%%%%%%%%%%%%%%%%%%%%%%%%%%%%%%%%%%%%%%%%%%%%%%%%%%%%%%%%%%%%%%%%%%%%%%

\begin{frame}{Discrete-time quantum walk on a line}

\begin{center}
\setlength{\unitlength}{.1cm}
\begin{picture}(100,5)(-10,-4)
%  \put(-10,-4){\framebox(100,5){}}
  \multiput(0,0)(10,0){9}{\circle*{1}}
  \thicklines
  \put(0,0){\line(1,0){80}}
  \put(0,0){\vector(-1,0){10}}
  \put(80,0){\vector(1,0){10}}
  \put(-1,-4){\makebox(2,3){-4}}
  \put(9,-4){\makebox(2,3){-3}}
  \put(19,-4){\makebox(2,3){-2}}
  \put(29,-4){\makebox(2,3){-1}}
  \put(39,-4){\makebox(2,3){0}}
  \put(49,-4){\makebox(2,3){1}}
  \put(59,-4){\makebox(2,3){2}}
  \put(69,-4){\makebox(2,3){3}}
  \put(79,-4){\makebox(2,3){4}}
\end{picture}
\end{center}

Add a ``coin'': state space $\spn\{|x\>\ot|{\leftarrow}\>,|x\>\ot|{\rightarrow}\>: x \in \Z\}$

\pause
\bigskip

Coin flip: $C := I \ot H$

\bigskip

Shift:
$\begin{array}{l@{\;}l}
S|x\>\ot|{\leftarrow}\> &= |x-1\>\ot|{\leftarrow}\> \\
S|x\>\ot|{\rightarrow}\> &= |x+1\>\ot|{\rightarrow}\>
\end{array}$

\bigskip

Walk step: $SC$

\pause
\begin{center}
  \includegraphics[width=7cm]{lineqd}
\end{center}

\end{frame}

%%%%%%%%%%%%%%%%%%%%%%%%%%%%%%%%%%%%%%%%%%%%%%%%%%%%%%%%%%%%%%%%%%%%%%%%%%%%%%%%

\begin{frame}{The Szegedy walk}

State space: $\spn\{|v\>\ot|w\>,|w\>\ot|v\>: (v,w) \in E\}$

\pause
\bigskip

Let $W$ be a stochastic matrix (a discrete-time random walk)

\begin{align*}
\onslide<3->{
\text{Define} \quad 
|\psi_v\> &:= |v\>\ot \sum_{w \in V} \sqrt{W_{wv}}|w\> \quad}
\onslide<4->{
\text{(note $\<\psi_v|\psi_w\> = \delta_{v,w}$)}} \\
\onslide<5->{R &:= 2 \sum_{v \in V}|\psi_v\>\<\psi_v| - I} \\
\onslide<6->{S(|v\>\ot|w\>) &:= |w\>\ot|v\>}
\end{align*}

\bigskip

\onslide<7->{
Then a step of the walk is the unitary operator $U := SR$
}

\end{frame}

%%%%%%%%%%%%%%%%%%%%%%%%%%%%%%%%%%%%%%%%%%%%%%%%%%%%%%%%%%%%%%%%%%%%%%%%%%%%%%%%

\begin{frame}{Spectrum of the walk}
  
Let $T := \sum_{v \in V} |\psi_v\>\<v|$, so $R=2TT^\dag - I$.

\pause
\bigskip
  
\begin{theorem}[Szegedy]
  Let $W$ be a stochastic matrix.
  Suppose the matrix
  \[
  \sum_{v,w} \sqrt{W_{vw}W_{wv}}|w\>\<v|
  \]
  has an eigenvector $|\lambda\>$ with eigenvalue $\lambda$.
  Then
  \[
    \frac{I - e^{\pm\ii\arccos\lambda}S}{\sqrt{2(1-\lambda^2)}} T|\lambda\>
  \]
  are eigenvectors of $U=SR$ with eigenvalues
  \[
    e^{\pm\ii\arccos\lambda}.
  \]
\end{theorem}

\end{frame}

%%%%%%%%%%%%%%%%%%%%%%%%%%%%%%%%%%%%%%%%%%%%%%%%%%%%%%%%%%%%%%%%%%%%%%%%%%%%%%%%

\begin{frame}{Proof of Szegedy's spectral theorem}

\begin{proof}[Proof sketch]
  Straightforward calculations give
  \[
    TT^\dag = \sum_{v \in V} |\psi_v\>\<\psi_v| \qquad\quad
    T^\dag T = I \]\[
    T^\dag S T = \sum_{v,w \in V} \sqrt{W_{vw}W_{wv}}|w\>\<v|
               = \sum_\lambda |\lambda\>\<\lambda|
  \]
  \pause
  which can be used to show
  \begin{align*}
    U(T|\lambda\>) &= ST|\lambda\> &
    U(ST|\lambda\>) &= 2\lambda ST|\lambda\> - T|\lambda\>.
  \end{align*}
  \pause
  Diagonalizing within the subspace $\spn\{T|\lambda\>,ST|\lambda\>\}$ gives the desired result.
\end{proof}

\pause
\bigskip

\lab{Exercise} Fill in the details

\end{frame}

%%%%%%%%%%%%%%%%%%%%%%%%%%%%%%%%%%%%%%%%%%%%%%%%%%%%%%%%%%%%%%%%%%%%%%%%%%%%%%%%

\begin{frame}{Random walk search algorithm}

Given $G=(V,E)$, let $M \subset V$ be a set of \emph{marked vertices}

\pause
\bigskip

Start at a random unmarked vertex

\bigskip

Walk until we reach a marked vertex\pause{}:
\begin{align*}
  W'_{vw} &:= \begin{cases}
  1 & w\in M \text{~and~} v=w \\
  0 & w\in M \text{~and~} v\ne w \\
  W_{vw} & w \notin M.
  \end{cases} \\
  &\onslide<4->{= \begin{pmatrix}
  W_M & 0 \\ V & I
\end{pmatrix}
\quad \text{($W_M$: delete marked rows and columns of $W$)}}
\end{align*}

\bigskip

\onslide<5->{
\lab{Question} How long does it take to reach a marked vertex?
}

\end{frame}

%%%%%%%%%%%%%%%%%%%%%%%%%%%%%%%%%%%%%%%%%%%%%%%%%%%%%%%%%%%%%%%%%%%%%%%%%%%%%%%%

\begin{frame}{Classical hitting time}
Take $t$ steps of the walk:
\begin{align*}
  (W')^t
  &=
  \begin{pmatrix}
  W_M^t & 0 \\ V (I + W_M + \cdots + W_M^{t-1}) & I
  \end{pmatrix} \\
  &\onslide<2->{= \begin{pmatrix}
  W_M^t & 0 \\ V \frac{I - W_M^t}{I - W_M} & I
  \end{pmatrix}}
\end{align*}

\medskip

\onslide<3->{
Convergence time depends on how close $\norm{W_M}$ is to $1$, which depends on the spectrum of $W$
}

\medskip

\begin{lemma}<4->
  Let $W = W^T$ be a symmetric Markov chain.
  Let the second largest eigenvalue of $W$ be $1-\delta$, and let $\epsilon = |M|/|V|$ (the fraction of marked items).
  Then the probability of reaching a marked vertex is $\Omega(1)$ after $t = O(1/\delta\epsilon)$ steps of the walk.
\end{lemma}

\end{frame}

%%%%%%%%%%%%%%%%%%%%%%%%%%%%%%%%%%%%%%%%%%%%%%%%%%%%%%%%%%%%%%%%%%%%%%%%%%%%%%%%

\begin{frame}{Quantum walk search algorithm}
  
Start from the state
$
%  \frac{1}{\sqrt{N-|M|}} \sum_{v \not\in M} T|v\> =
  \frac{1}{\sqrt{N-|M|}} \sum_{v \not\in M} |\psi_v\>
$

\pause
\bigskip

Consider the walk $U$ corresponding to $W'$\pause{}:
\[
  \sum_{v,w \in V} \sqrt{W'_{v,w}W'_{w,v}}|w\>\<v|
  = \begin{pmatrix}
  W_M & 0 \\ 0 & I
\end{pmatrix}
\] \pause
Eigenvalues of $U$ are $e^{\pm\ii\arccos\lambda}$ where the $\lambda$ are eigenvalues of $W_M$

\pause
\bigskip

Perform phase estimation on $U$ with precision $O(\sqrt{\delta\epsilon})$
\begin{itemize}
  \item<+-> no marked items $\implies$ estimated phase is $0$
  \item<+-> $\epsilon$ fraction of marked items $\implies$ nonzero phase with probability $\Omega(1)$
\end{itemize}

\bigskip

\onslide<7->
Further refinements give algorithms for \emph{finding} a marked item

\end{frame}

%%%%%%%%%%%%%%%%%%%%%%%%%%%%%%%%%%%%%%%%%%%%%%%%%%%%%%%%%%%%%%%%%%%%%%%%%%%%%%%%

\begin{frame}{Grover's algorithm revisited}
  
\begin{problem}
Given a black box $f:X \to \{0,1\}$, is there an $x$ with $f(x)=1$?
\end{problem}

\pause
\bigskip

Markov chain on $N=|X|$ vertices:
\[
  W := \frac{1}{N} \begin{pmatrix}
  1 & \cdots & 1 \\
  \vdots & \ddots & \vdots \\
  1 & \cdots & 1
  \end{pmatrix}
  = |\psi\>\<\psi|,
  \quad
  |\psi\> := \frac{1}{\sqrt N}\sum_{x \in X}|x\>
\] \pause
Eigenvalues of $W$ are $0,1$ $\implies$ $\delta = 1$

\pause
\bigskip

Hard case: one marked vertex, $\epsilon = 1/N$
\pause
\bigskip

Hitting times
\begin{itemize}
  \item<+-> Classical: $O(1/\delta\epsilon) = O(N)$
  \item<+-> Quantum: $O(1/\sqrt{\delta\epsilon}) = O(\sqrt{N})$
\end{itemize}

\end{frame}

%%%%%%%%%%%%%%%%%%%%%%%%%%%%%%%%%%%%%%%%%%%%%%%%%%%%%%%%%%%%%%%%%%%%%%%%%%%%%%%%

\begin{frame}{Element distinctness}

\begin{problem}
Given a black box $f:X \to Y$, are there distinct $x,x'$ with $f(x)=f(x')$?
\end{problem}

\pause
\bigskip

Let $N=|X|$; classical query complexity is $\Omega(N)$

\pause
\bigskip

Consider a quantum walk on the Hamming graph $H(N,M)$ \pause
\begin{itemize}
  \item<+-> Vertices: $\{(x_1,\ldots,x_M): x_i \in X\}$
  \item<+-> Store the values $(f(x_1),\ldots,f(x_M))$ at vertex $(x_1,\ldots,x_M)$
  \item<+-> Edges between vertices that differ in exactly one coordinate
\end{itemize}

\end{frame}

%%%%%%%%%%%%%%%%%%%%%%%%%%%%%%%%%%%%%%%%%%%%%%%%%%%%%%%%%%%%%%%%%%%%%%%%%%%%%%%%

\begin{frame}{Element distinctness: Analysis}
  
Spectral gap: $\delta = O(1/M)$

\pause
\bigskip

Fraction of marked vertices: $\epsilon \ge \binom{N-2}{M-2}/\binom{N}{M} = \Theta(M^2/N^2)$

\pause
\bigskip

Quantum hitting time: $O(1/\sqrt{\delta\epsilon}) = O(N/\sqrt M)$

\pause
\bigskip

Quantum query complexity: \pause
\begin{itemize}
  \item<+-> $M$ queries to prepare the initial state
  \item<+-> $2$ queries for each step of the walk (compute $f$, uncompute $f$)
  \item<+-> Overall: $M + O(N/\sqrt M)$
\end{itemize}

\bigskip

\onslide<8->
Choose $M=N^{2/3}$: query complexity is $O(N^{2/3})$ ~~~ (optimal!)

\end{frame}

%%%%%%%%%%%%%%%%%%%%%%%%%%%%%%%%%%%%%%%%%%%%%%%%%%%%%%%%%%%%%%%%%%%%%%%%%%%%%%%%

\begin{frame}{Quantum walk algorithms}

Quantum walk search algorithms
\begin{itemize}
  \item Spatial search
  \item Finding a triangle in a graph
  \item Checking matrix multiplication
  \item Testing if a black-box group is abelian
\end{itemize}

\bigskip

Evaluating Boolean formulas

\bigskip

Exponential speedup for a natural problem?

\end{frame}